\section{Tipos de Mantenimiento}

Una vez realizada la lista de equipos, desglosados incluso en los elementos que los componen e identificado cada item con un código único que
permite referenciarlo, la siguiente tarea que debemos abordar es la de decidir
cómo vamos a mantener cada uno de esos equipos.
Tradicionalmente, se han distinguido 5 tipos de mantenimiento, que se diferencian entre sí por el carácter de las tareas que incluyen:

\begin{tcolorbox}[]
DIVISIÓN CLÁSICA DE TIPOS DE MANTENIMIENTO
\tcblower
\begin{itemize}
\item Mantenimiento correctivo.
\item Mantenimiento preventivo.
\item Mantenimiento predictivo.
\item Mantenimiento hard time o cero horas.
\item Mantenimiento en uso.
\end{itemize}
\end{tcolorbox}

\begin{itemize}
\item Mantenimiento correctivo: Es el conjunto de tareas destinadas a corregir los defectos que se van presentando en los distintos equipos y
que son comunicados al departamento de mantenimiento por los
usuarios de los mismos.
\item Mantenimiento preventivo: Es el mantenimiento que tiene por misión
mantener un nivel de servicio determinado en los equipos, programando las correcciones de sus puntos vulnerables en el momento
más oportuno.
\item Mantenimiento predictivo: Es el que persigue conocer e informar permanentemente del estado y operatividad de las instalaciones mediante
el conocimiento de los valores de determinadas variables, representativas de tal estado y operatividad. Para aplicar este mantenimiento es
necesario identificar variables físicas (temperatura, vibración, consumo
de energía, etc.) cuya variación sea indicativa de problemas que puedan
estar apareciendo en el equipo. Es el tipo de mantenimiento más tecnológico, pues requiere de medios técnicos avanzados, y de fuertes conocimientos matemáticos, físicos y técnicos.
\item Mantenimiento cero horas: Es el conjunto de tareas cuyo objetivo es revisar los equipos a intervalos programados bien antes de que aparezca
ningún fallo, bien cuando la fiabilidad del equipo ha disminuido apreciablemente, de manera que resulta arriesgado hacer previsiones sobre
su capacidad productiva. Dicha revisión consiste en dejar el equipo a
cero horas de funcionamiento, es decir, como si el equipo fuera nuevo.
En estas revisiones se sustituyen o se reparan todos los elementos sometidos a desgaste. Se pretende asegurar, con gran probabilidad, un
tiempo de buen funcionamiento fijado de antemano.
\item Mantenimiento en uso: es el mantenimiento básico de un equipo realizado por los usuarios del mismo. Consiste en una serie de tareas elementales (tomas de datos, inspecciones visuales, limpieza, lubricación, reapriete de tornillos) para las que no es necesario una gran
formación, sino tan solo un entrenamiento breve. Este tipo de mantenimiento es la base del TPM (Total Productive Maintenance, Mantenimiento Productivo Total)
\end{itemize}

\subsubsection*{LOS TIPOS DE MANTENIMIENTO NO SON
DIRECTAMENTE APLICABLES}
Esta división de Tipos de Mantenimiento presenta el inconveniente de que
cada equipo necesita una mezcla de cada uno de esos tipos, de manera que no
podemos pensar en aplicar uno solo de ellos a un equipo en particular.
Así, en un motor determinado nos ocuparemos de su lubricación (mantenimiento preventivo periódico), si lo requiere, mediremos sus vibraciones o
sus temperaturas (mantenimiento predictivo), quizás le hagamos una puesta a
punto anual (puesta a cero) y repararemos las averías que vayan surgiendo
(mantenimiento correctivo). La mezcla más idónea de todos estos tipos de
mantenimiento nos la dictarán estrictas razones ligadas al coste de las pérdidas de producción en una parada de ese equipo, al coste de reparación, al impacto ambiental, a la seguridad y a la calidad del producto o servicio, entre
otras.
El inconveniente, pues, de la división anterior es que no es capaz de dar
una respuesta clara a esta pregunta:

¿Cuál es el mantenimiento que debo aplicar a cada uno de los equipos que componen una planta concreta?

Para dar respuesta a esta pregunta es conveniente definir el concepto de modelo de mantenimiento. Un modelo de mantenimiento es una mezcla de los anteriores tipos de mantenimiento en unas proporciones determinadas, y que
responde adecuadamente a las necesidades de un equipo concreto. Podemos
pensar que cada equipo necesitará una mezcla distinta de los diferentes tipos
de mantenimiento, una mezcla determinada de tareas, de manera que los
modelos de mantenimiento posibles serán tantos como equipos puedan existir. Pero esto no es del todo correcto. Pueden identificarse claramente 4 de estas mezclas, complementadas con otros dos tipos de tareas adicionales, según
veremos más adelante.

\actividad{Leer,contestar e investigar}
\begin{enumerate}
    \item ¿Qué es T.P.M. y en que consiste?
    \item ¿Cuál es la diferencia entre modelos de mantenimiento y tipo de mantenimiento?
    \item El mantenimiento predictivo involucra conocimientos técnicos de diversas herramientas tecnológicas, investigar brevemente cada una:
    \begin{itemize}
        \item Rayos X.
        \item Ultrasonido.
        \item Termografía
        \item Vibraciones.
        \item Tribología.
        \item Ensayos de Aislamiento eléctrico.
        \item Liquidos Penetrantes y Partículas Magnéticas.
    \end{itemize}
\end{enumerate}
    